% 英文摘要

With the popularization of online education, platforms like MOOCs have provided learners with abundant learning resources. However, they also face challenges such as large individual differences among learners, difficulty in providing personalized guidance, and a lack of timely and effective feedback. Traditional online learning platforms mostly display static content, which makes it hard to adapt to the backgrounds and paces of different learners. In addition, in core scenarios such as code training and video learning, students often cannot receive targeted heuristic tutoring when they encounter problems, leading to low learning efficiency. To solve the above problems, this thesis designs and implements a personalized learning system based on Large Language Model (LLM) and Retrieval-Augmented Generation (RAG) technologies.

This thesis first conducts a comprehensive analysis of the system requirements, specifying core functions such as knowledge base-driven video Q\&A, intelligent code training, and code analysis based on Abstract Syntax Tree (AST). In terms of system architecture, a front-end and back-end separation pattern is adopted. The front-end uses the Vue 3 framework, the back-end is built on Flask, and MySQL and Redis are utilized for data and state management. On this basis, the thesis deeply discusses the implementation of key technologies in the system. For the code training module, an algorithm combining Socratic heuristic strategies is proposed and implemented. By parsing the AST structure of the code submitted by students and combining it with the reasoning ability of the LLM, the algorithm provides students with progressive debugging guidance instead of direct answers, thereby cultivating their independent problem-solving skills. In the video learning module, the system introduces RAG technology to convert the subtitles and documents of course videos into vectors and store them in the FAISS database. When learners ask questions, the system can retrieve relevant contexts within the knowledge base of a specific video and generate answers with precise timestamp citations, enabling accurate Q\&A interactions.

Finally, a comprehensive functional and performance test of the system is conducted. The test results show that the system not only achieves 100\% code coverage in white-box testing but also performs stably in concurrent performance testing, meeting the usage requirements in actual teaching scenarios. The design and implementation of this system not only verify the application potential of large model technologies in the vertical field of education but also provide useful exploration and practical references for the future construction of intelligent and personalized online education platforms.